\section{Introduction}
\label{sec:intro}

Software is a state machine, and Part I of this work~\citep{welcing2026nerve} asked what it takes to ship a small one as a network rather than as code: a network that takes the interface's
state and a user's action and returns the next state, held to a behaviour specification the way code is held
to types and tests. The answer was a spec language with a derived oracle, a compiler from the explicit part of
a spec to exact weights, an energy model for the rest, and a protocol under which seven small finite interfaces were
learned-only exact on cells they had never seen and a continuous slider was within $1\%$ of range. It was a result about one interface at a time, and its own
limitations section said so: nothing composed, the largest interface had $404$ states, and the only inputs
were a click and a drag.

Two things have changed since. The first is the goal. An interface's network is trained on that interface, on
the whole of its spec, so whether it would have generalised from part of the spec is beside the point; the
questions are a compiler's: exact, every time, at what cost. The second is who else is compiling machines into
weights. Qwenputer and qwenGPU~\citep{qwenputer,qwengpu} compile a 32-bit RISC-V CPU and a SIMT GPU into layers
in front of a language model of two to four billion parameters, with the model's own recurrent state as the
machine's memory, so that a general machine lives inside a model that drives it through its own tool calls.
That is one road to programs in weights, and it is not ours. Our parts are four orders of magnitude smaller,
cannot talk, and are exact by a check rather than by a construction; one of our steps is one network
evaluation, where one word of a program's output from Qwenputer is tens of thousands of instruction tokens and
tens of seconds. The question our parts raise is how they become something as large as a desktop application
without becoming one large network, which the training cost forbids.

This paper is in two halves. The first (\cref{sec:standing}) is measured: of the \nPlanned{} interfaces we set out to
train, \nNewDemos{} of them rebuilt from the demo set of declarative dragging~\citep{dragology} because that
library was built to show how much a drag can mean, one network per interface trained on its whole spec is
exact on every finite interface that has finished (\nExactSpecs{} of \nFiniteSpecs{} as this draft is built),
on every seed, with no rule broken, including a Tower of Hanoi whose solver the spec shows only by
example; restricting answers to the states the invariants allow removes rule breaks on a starved network without
retraining; and the recipe's cost grows with states times cells until it stops. The second half
(\cref{sec:blocks,sec:arch,sec:runtime,sec:demos,sec:protocol}) is a plan, written as a paper so that it can be
scored. We propose that large interactive systems be \emph{machines}: blocks that own their state and expose
ports; three couplings, independent, event-wired and read-wired, each with a training cost; an argument that
exactness composes under correct wiring; a block architecture whose heads are chosen by the types of a spec's
fields and whose candidates come from a proposer when the state space is large; a runtime with a synchronous
tick, bindings for keys, modifiers and the mouse, selection, focus and undo, drags inside a running machine, and
a rule that keeps continuous motion exact across hand-offs; and five demos, from an odometer of digit networks
to a desktop card board, each registered with a prediction and the result that would falsify it. The picture
the plan is built toward is literal: a machine you can watch run through a series of separate networks, each
hand-off checked against the spec as it happens, that you can reach into and drag while it runs.

\paragraph{Contributions.}
(1)~One network per interface, trained on its whole spec, is exact: \nPlanned{} interfaces, \nNewDemos{} from
Dragology's demo set, with generated tables, scripted provenance, and a valid-candidate rule that removes rule
breaks without retraining (\cref{sec:standing}).
(2)~A calculus of blocks, ports and wiring with a cost model for three couplings and a proof sketch that
exactness composes (\cref{sec:blocks}).
(3)~A block architecture: heads by field type, a proposer for large spaces, a compiled core with trained holes,
the energy as a port type, and explicit size and latency budgets (\cref{sec:arch}).
(4)~A runtime: the synchronous tick, a declarative input model, selection, focus and undo, drags inside running
machines, discrete hand-offs between continuous blocks, and derived rendering (\cref{sec:runtime}).
(5)~Five demos registered as experiments with predictions and falsifiers, and an evaluation protocol for
machines of networks (\cref{sec:demos,sec:protocol}).
